
\documentclass{article}
\usepackage{colortbl}
\usepackage{makecell}
\usepackage{multirow}
\usepackage{supertabular}

\begin{document}

\newcounter{utterance}

\centering \large Interaction Transcript for game `cladder', experiment `full\_v1.5\_default', episode 8348 with llama{-}3.1{-}8b{-}instruct.
\vspace{24pt}

{ \footnotesize  \setcounter{utterance}{1}
\setlength{\tabcolsep}{0pt}
\begin{supertabular}{c@{$\;$}|p{.15\linewidth}@{}p{.15\linewidth}p{.15\linewidth}p{.15\linewidth}p{.15\linewidth}p{.15\linewidth}}
    \# & \multicolumn{2}{c}{Player} && \multicolumn{2}{c}{Game Master} \\
    \hline

    \theutterance \stepcounter{utterance}  
    & & & \multicolumn{4}{p{0.6\linewidth}}{
        \cellcolor[rgb]{0.9,0.9,0.9}{
            \makecell[{{p{\linewidth}}}]{
                \texttt{\tiny{[P1$\langle$GM]}}
                \texttt{You are an expert in causal inference. The following question is not a typical commonsense query, but rather a meticulously designed question created by a professor specializing in causal inference, intended to assess the students' mastery of the course content.} \\
\texttt{Imagine a self{-}contained, hypothetical world with only the following conditions, and without any unmentioned factors or causal relationships: Yomx has a direct effect on cwoi and gwet. Gwet has a direct effect on xevu. Cwoi has a direct effect on xevu. The overall probability of yomx is 22\%. The probability of not yomx and xevu is 48\%. The probability of yomx and xevu is 7\%. Is the chance of xevu larger when observing yomx?} \\
\texttt{Start your answer with 'Yes' or 'No', followed by additional reasoning or evidence to support your explanation.} \\
            }
        }
    }
     \\ \\

    \theutterance \stepcounter{utterance}  
    & \multicolumn{4}{p{0.6\linewidth}}{
        \cellcolor[rgb]{0.561,0.737,0.561}{
            \makecell[{{p{\linewidth}}}]{
                \texttt{\tiny{[P1$\rangle$GM]}}
                \texttt{No, the chance of xevu is not larger when observing yomx.} \\
\\ 
\texttt{To determine the causal effect of yomx on xevu, we can use the backdoor criterion, which involves identifying a path from yomx to xevu that is unconfounded by any common causes. } \\
\\ 
\texttt{We have two paths from yomx to xevu: } \\
\\ 
\texttt{1. yomx {-}> cwoi {-}> xevu} \\
\texttt{2. yomx {-}> gwet {-}> xevu} \\
\\ 
\texttt{However, we also have a direct effect of cwoi on xevu, which creates a confounding variable. To resolve this, we can use the backdoor adjustment formula, which involves calculating the probability of xevu given yomx and the probability of xevu given not yomx, and then combining these probabilities using the formula:} \\
\\ 
\texttt{P(xevu|yomx) = P(xevu|yomx) + P(xevu|not yomx) {-} P(xevu|not yomx and yomx)} \\
\\ 
\texttt{We are given the following probabilities:} \\
\\ 
\texttt{{-} P(yomx) = 0.22} \\
\texttt{{-} P(not yomx and xevu) = 0.48} \\
\texttt{{-} P(yomx and xevu) = 0.07} \\
\\ 
\texttt{We can calculate P(xevu|not yomx) as follows:} \\
\\ 
\texttt{P(xevu|not yomx) = P(not yomx and xevu) / P(not yomx)} \\
\texttt{= 0.48 / (1 {-} 0.22)} \\
\texttt{= 0.48 / 0.78} \\
\texttt{= 0.615} \\
\\ 
\texttt{We can also calculate P(xevu|not yomx and yomx) as follows:} \\
\\ 
\texttt{P(xevu|not yomx and yomx) = P(not yomx and yomx and xevu) / P(not yomx and yomx)} \\
\texttt{= 0 / (1 {-} 0.22) * 0.22} \\
\texttt{= 0 / 0.78 * 0.22} \\
\texttt{= 0} \\
\\ 
\texttt{Now, we can plug in the values into the backdoor adjustment formula:} \\
\\ 
\texttt{P(xevu|yomx) = P(yomx and xevu) / P(yomx) + P(xevu|not yomx) {-} P(xevu|not yomx and yomx)} \\
\texttt{= 0.07 / 0.22 + 0.615 {-} 0} \\
\texttt{= 0.3182 + 0.615} \\
\texttt{= 0.9332} \\
\\ 
\texttt{P(xevu|not yomx) = 0.615} \\
\\ 
\texttt{Since P(xevu|yomx) = 0.9332 is greater than P(xevu|not yomx) = 0.615,} \\
            }
        }
    }
    & & \\ \\

    \theutterance \stepcounter{utterance}  
    & & & \multicolumn{2}{p{0.3\linewidth}}{
        \cellcolor[rgb]{0.9,0.9,0.9}{
            \makecell[{{p{\linewidth}}}]{
                \texttt{\tiny{[GM$|$GM]}}
                \texttt{No} \\
            }
        }
    }
    & & \\ \\

    \theutterance \stepcounter{utterance}  
    & & & \multicolumn{2}{p{0.3\linewidth}}{
        \cellcolor[rgb]{0.9,0.9,0.9}{
            \makecell[{{p{\linewidth}}}]{
                \texttt{\tiny{[GM$|$GM]}}
                \texttt{no} \\
            }
        }
    }
    & & \\ \\

    \theutterance \stepcounter{utterance}  
    & & & \multicolumn{2}{p{0.3\linewidth}}{
        \cellcolor[rgb]{0.9,0.9,0.9}{
            \makecell[{{p{\linewidth}}}]{
                \texttt{\tiny{[GM$|$GM]}}
                \texttt{game\_result = WIN} \\
            }
        }
    }
    & & \\ \\

\end{supertabular}
}

\end{document}
